\chapter{Mathematical Proofs}
\label{chp:appB}

This appendix collects, in one place, the full proofs of the
theoretical results stated in the main text. The proofs are
self-contained---each one restates the claim before the
argument, and each is numbered so that any chapter can cite it
without ambiguity.

\section{Notation and Conventions}\label{app:notation}
The following conventions are used throughout this appendix.
\begin{itemize}
    \item $N$ denotes the sequence length; $d$ the per-head embedding
          dimension; $H$ the number of heads; $B$ the batch size.
    \item $(B_r, B_c)$ are the row/column block sizes of Flash
          Attention; $T_r = \lceil N/B_r \rceil$ and
          $T_c = \lceil N/B_c \rceil$.
    \item $(w_l, w_r)$ are the left/right window half-widths; the full
          window is $w = w_l + w_r + 1$.
    \item $Q, K \in \mathbb{R}^{N \times d}$ and $V \in \mathbb{R}^{N
          \times d_v}$ are the query/key/value matrices.
    \item $A_{ij} \triangleq q_i \cdot k_j / \sqrt{d}$ is the raw
          attention score; the softmax-normalised weight is $P_{ij}$.
\end{itemize}

\section{Correctness of Corner-Based Block
Classification}\label{app:proofs_classify}

\begin{lemma}[Restatement of Lemma~\ref{lem:corner}]\label{lem:corner2}
Let $B = [r_0, r_1] \times [c_0, c_1]$ be a block of the attention
matrix. Define the window band
$W = \{(i,j) \in [0, N-1]^2 : -w_l \le j - i \le w_r\}$.
Then:
\begin{enumerate}
    \item[\textup{(i)}] $B \cap W = \emptyset$ iff
          $c_1 < r_0 - w_l$ \textbf{or} $c_0 > r_1 + w_r$;
    \item[\textup{(ii)}] $B \subseteq W$ iff
          $c_1 \le r_0 + w_r$ \textbf{and} $c_0 \ge r_1 - w_l$.
\end{enumerate}
\end{lemma}

\begin{proof}
Define the scalar function $\phi(i, j) = j - i$ on the block
$B$. Since $\phi$ is linear and $B$ is an axis-aligned
rectangle, $\phi$ attains its minimum at the corner $(r_1, c_0)$
(the corner with the largest $i$ and the smallest $j$) and its
maximum at $(r_0, c_1)$ (smallest $i$, largest $j$).

\textbf{(i) Outside case.} $B \cap W = \emptyset$ holds iff
$\phi$ never lands in $[-w_l, w_r]$ on $B$, which is equivalent
to requiring $\max_B \phi < -w_l$ \emph{or} $\min_B \phi > w_r$.
Plugging in $\max_B \phi = c_1 - r_0$ and
$\min_B \phi = c_0 - r_1$ gives
$c_1 - r_0 < -w_l$ or $c_0 - r_1 > w_r$, i.e.\
$c_1 < r_0 - w_l$ or $c_0 > r_1 + w_r$.

\textbf{(ii) Inside case.} $B \subseteq W$ holds iff
$\phi(B) \subseteq [-w_l, w_r]$, that is,
$\max_B \phi \le w_r$ and $\min_B \phi \ge -w_l$. Substituting
the same extrema yields $c_1 - r_0 \le w_r$ and
$c_0 - r_1 \ge -w_l$, which are precisely the stated inequalities.

Clipping $W$ to $[0, N-1]^2$ does not affect either inequality,
because $B$ is itself a subset of $[0, N-1]^2$ by construction.
\end{proof}

\begin{corollary}\label{cor:classifier_cost}
Algorithm~\ref{alg:swfa_classify} runs in $\mathcal{O}(T_r T_c)$
arithmetic operations and correctly assigns each block to exactly one
of the sets $\mathcal{I}, \mathcal{X}, \mathcal{O}$.
\end{corollary}

\begin{proof}
For each of the $T_r T_c$ blocks,
Algorithm~\ref{alg:swfa_classify} performs a constant number of
comparisons derived from the two predicates of
Lemma~\ref{lem:corner2}. Disjointness of the three sets is
enforced by the order of the tests: the \texttt{outside}
predicate takes precedence; the \texttt{inside} predicate is
tested only when \texttt{outside} is false; any remaining block
is assigned to $\mathcal{X}$. By Lemma~\ref{lem:corner2}, a
block is outside exactly when \texttt{outside} is true and
inside exactly when \texttt{inside} is true. The three sets
therefore form a partition of the grid.
\end{proof}

\section{Block-Count Lemma}\label{app:proofs_block}

\begin{lemma}[Skipped-block count]\label{lem:blockcount}
For a symmetric window $w_l = w_r = W$ and equal block sizes
$B_r = B_c = B$, with $B \mid N$, the number of outside blocks is
\begin{equation}\label{eq:block_count}
    |\mathcal{O}|
    \;=\; 2 \sum_{k=1}^{T_r - 1}
          \max(0, T_r - k - \lceil (W + (k-1) B) / B \rceil)
    \;=\; (T_r)^2 - \Theta(T_r W / B)
\end{equation}
and therefore the skip ratio converges to $1 - 2W / N$ as
$N \to \infty$ with fixed $W$.
\end{lemma}

\begin{proof}
Under the assumption $B \mid N$, the grid contains $T = N/B$
blocks along each axis. Lemma~\ref{lem:corner2}(i) then says that
block $(r, c)$ is outside iff
$c B + B - 1 < r B - W$ or $c B > r B + B - 1 + W$, which
simplifies to
\begin{equation}\label{eq:blockout}
    c \le r - \left\lceil (W + B - 1)/B \right\rceil
    \quad \text{or} \quad
    c \ge r + \left\lceil (W + B - 1)/B \right\rceil.
\end{equation}
Writing $k^\star = \lceil (W + B - 1)/B \rceil$, the number of
outside blocks in row $r$ becomes
$\max(0, r - k^\star) + \max(0, T - 1 - r - k^\star)$. Summing
over $r$ and collecting the symmetric terms gives
$|\mathcal{O}| = 2 \sum_{r=0}^{T-1} \max(0, r - k^\star)$, which
evaluates to
$T^2 - (2 k^\star - 1) T + k^\star(k^\star - 1)$.
Dividing through by $T^2$ and taking $N \to \infty$ at fixed $W$
gives $k^\star / T \to W / N$, so the skip ratio converges to
$1 - 2W/N$, as claimed.
\end{proof}

\section{Convexity of the SWFA Cost}\label{app:proof_convex}

\begin{theorem}[Restatement of
Theorem~\ref{thm:swfa_convex}]\label{thm:convex2}
The cost function
\begin{equation}
    \mathcal{L}(w) = \alpha \kappa N w d + \beta M_0
        + \frac{2 \gamma N A_0}{1 - e^{-\lambda}}\, e^{-\lambda w / 2}
\end{equation}
is strictly convex for $w > 0$.
\end{theorem}

\begin{proof}
Differentiating $\mathcal{L}$ twice yields
\begin{equation}
    \mathcal{L}'(w) = \alpha \kappa N d
      - \frac{\gamma N A_0 \lambda}{1 - e^{-\lambda}}
        e^{-\lambda w / 2},
\end{equation}
\begin{equation}
    \mathcal{L}''(w)
      = \frac{\gamma N A_0 \lambda^2}{2(1 - e^{-\lambda})}
        e^{-\lambda w / 2}.
\end{equation}
Every prefactor is positive: $\gamma > 0$ is assumed,
$N$, $A_0$, and $\lambda$ are positive by the model
specification, and $1 - e^{-\lambda} > 0$ whenever $\lambda > 0$.
The exponential $e^{-\lambda w / 2}$ is strictly positive for
every $w \in \mathbb{R}$. Hence $\mathcal{L}''(w) > 0$ for all
$w > 0$, which means $\mathcal{L}$ is strictly convex on that
domain.
\end{proof}

\begin{remark}
$\mathcal{L}''$ is a positive exponential multiplied by a
constant, so the convexity is in fact \emph{strong}: for any
$w_0 > 0$ there exists $\mu > 0$ such that
$\mathcal{L}''(w) \ge \mu$ on any bounded interval containing
$w_0$. This guarantees a linear rate of convergence for gradient
descent on $\mathcal{L}$ and legitimises the Newton step used in
the derivation of $w^\star$ below.
\end{remark}

\section{Closed-Form Optimal Window}\label{app:proof_optw}

\begin{theorem}[Restatement of
Theorem~\ref{thm:swfa_optw}]\label{thm:optw2}
The minimiser of $\mathcal{L}$ is
\begin{equation}\label{eq:opt_w_app}
    w^\star = \frac{2}{\lambda}
      \ln\!\left(
      \frac{\gamma A_0 \lambda}{\alpha \kappa d (1 - e^{-\lambda})}
      \right),
\end{equation}
provided the argument of the logarithm exceeds one (otherwise the
unconstrained minimiser is non-positive and the projected minimiser
equals $1$).
\end{theorem}

\begin{proof}
By Theorem~\ref{thm:convex2} the loss $\mathcal{L}$ is strictly
convex on the domain $w > 0$, so its minimiser coincides with
its unique stationary point. Setting $\mathcal{L}'(w) = 0$ gives
\begin{equation}
    \alpha \kappa N d
    = \frac{\gamma N A_0 \lambda}{1 - e^{-\lambda}}
      e^{-\lambda w / 2},
\end{equation}
and therefore
\begin{equation}
    e^{-\lambda w / 2}
    = \frac{\alpha \kappa d (1 - e^{-\lambda})}{\gamma A_0 \lambda}.
\end{equation}
Applying $\ln$ to both sides and multiplying by $-2/\lambda$
recovers~\eqref{eq:opt_w_app}. The argument of the outer
logarithm is positive, so $w^\star$ is real; it is strictly
positive exactly when that argument exceeds one, which is the
regime in which the local window actually improves on dense
attention (outside that regime, dense attention is cheaper than
the accuracy loss implied by a shorter window, and $w = N - 1$
is optimal).

Projection onto $[1, N-1]$ preserves uniqueness. At the left
endpoint, convexity gives $\mathcal{L}(1) \le \mathcal{L}(0)$;
at the right endpoint, $\mathcal{L}(N-1)$ is dominated by the
computation term because the exponential has already decayed.
The projected minimiser is thus obtained by clipping $w^\star$
to $[1, N-1]$.
\end{proof}

\begin{corollary}[Sensitivities]\label{cor:sens}
Let $w^\star$ be as in Theorem~\ref{thm:optw2}. Then:
\begin{enumerate}
    \item[\textup{(i)}] $\partial w^\star / \partial \alpha < 0$;
    \item[\textup{(ii)}] $\partial w^\star / \partial \gamma > 0$;
    \item[\textup{(iii)}] $\partial w^\star / \partial d < 0$;
    \item[\textup{(iv)}] $w^\star \to \infty$ as $\lambda \to 0^+$.
\end{enumerate}
\end{corollary}

\begin{proof}
Taking $\partial / \partial \alpha$
of~\eqref{eq:opt_w_app} gives
$\partial w^\star / \partial \alpha = -2/(\lambda \alpha) < 0$.
Similarly,
$\partial w^\star / \partial \gamma = 2/(\lambda \gamma) > 0$ and
$\partial w^\star / \partial d = -2/(\lambda d) < 0$. For (iv), expand
the logarithm for small $\lambda$: $1 - e^{-\lambda}
\approx \lambda - \lambda^2/2$, so the argument of the outer log tends
to $\gamma A_0 / (\alpha \kappa d)$, which is $O(1)$. The outer factor
$2/\lambda$ diverges, so $w^\star \to \infty$.
\end{proof}

\section{Flash Attention Softmax
Correctness}\label{app:proof_flash_softmax}

Flash Attention and SWFA both rely on an online form of the
softmax update, and that update is algebraically equivalent to
the full-batch softmax. The argument is a short induction and is
reproduced below for completeness.

\begin{lemma}[Online softmax equivalence]\label{lem:online_softmax}
Given scalars $s_1, \dots, s_N$ and values $v_1, \dots, v_N$, let
$m = \max_i s_i$,
$\ell = \sum_i e^{s_i - m}$, and $o = \sum_i e^{s_i - m} v_i / \ell$.
Then for any partition of the indices into $T$ blocks,
the iterative update
\begin{align}
    m^{(t)} &= \max(m^{(t-1)}, \max_{i \in B_t} s_i),\\
    \ell^{(t)} &= e^{m^{(t-1)} - m^{(t)}} \ell^{(t-1)}
                 + \sum_{i \in B_t} e^{s_i - m^{(t)}},\\
    o^{(t)} &= \frac{e^{m^{(t-1)} - m^{(t)}} \ell^{(t-1)} o^{(t-1)}
                      + \sum_{i \in B_t} e^{s_i - m^{(t)}} v_i}
                    {\ell^{(t)}},
\end{align}
with $(m^{(0)}, \ell^{(0)}, o^{(0)}) = (-\infty, 0, 0)$, yields
$(m^{(T)}, \ell^{(T)}, o^{(T)}) = (m, \ell, o)$.
\end{lemma}

\begin{proof}
We argue by induction on $t$. Let $I_t = \bigcup_{s \le t} B_s$
denote the set of indices processed by step $t$, and define
$m_t = \max_{i \in I_t} s_i$,
$\ell_t = \sum_{i \in I_t} e^{s_i - m_t}$, and
$o_t = \sum_{i \in I_t} e^{s_i - m_t} v_i / \ell_t$.

\emph{Base case.} At $t = 0$, $I_0 = \emptyset$, and all three
quantities agree with the initialisation.

\emph{Inductive step.} Suppose the claim holds at step $t-1$. At
step $t$, $m^{(t)} = \max(m^{(t-1)}, \max_{i \in B_t} s_i) = m_t$
by definition. Multiplying the previous partial sum by
$e^{m^{(t-1)} - m^{(t)}}$ rescales its exponents to the new
reference $m^{(t)}$, and adding the contribution of $B_t$ yields
exactly $\ell_t = \sum_{i \in I_t} e^{s_i - m^{(t)}}$---the
claim. Applying the same rescaling to the numerator of $o$
reproduces $o_t$.
\end{proof}

\begin{corollary}[SWFA equivalence]
SWFA produces the same output as vanilla sliding-window attention up
to floating-point round-off.
\end{corollary}

\begin{proof}
Lemma~\ref{lem:corner2}(i) guarantees that SWFA skips precisely
those blocks that lie outside the window. On the blocks that
remain, SWFA either applies no mask (inside blocks) or applies
the additive $\{0, -\infty\}$ mask of
Algorithm~\ref{alg:swfa_mask} (intersecting blocks). Masked
entries contribute zero to both the numerator and the
denominator of the softmax, so the online update of
Lemma~\ref{lem:online_softmax}, applied to the masked scores,
is equivalent to the same update applied to the full collection
of in-window scores. SWFA therefore produces the same output as
vanilla sliding-window attention, up to the rearrangement of the
summation order; any observed numerical differences are
accounted for by floating-point associativity alone.
\end{proof}

\section{Communication-Cost Bounds for
Distributed Strategies}\label{app:proof_comms}

\begin{lemma}[DDP ring-allreduce bandwidth]\label{lem:ddp_bw}
On a ring topology with $P$ workers and link bandwidth $B$, a single
all-reduce of a tensor of size $M s$ bytes takes
\begin{equation}\label{eq:ring_cost}
    T_{\text{ring}}(M, P) =
    2 (P - 1) \alpha + \frac{2 (P - 1) M s}{P B},
\end{equation}
where $\alpha$ is the per-hop latency and the factor $2(P-1)/P$
reflects the two phases (reduce-scatter + all-gather).
\end{lemma}

\begin{proof}
A ring all-reduce runs in two phases of $P - 1$ steps each, and
every step forwards a fragment of size $M s / P$ to the next
worker. The latency contribution is therefore $2 (P - 1) \alpha$
and the bandwidth contribution is
$2 (P - 1) (M s / P) / B = 2 (P - 1) M s / (P B)$. Adding the
two gives~\eqref{eq:ring_cost}.
\end{proof}

\begin{corollary}\label{cor:ddp_limit}
Letting $P \to \infty$ with $M$ held fixed,
$T_{\text{ring}}(M, P) \to 2 M s / B + 2 (P - 1) \alpha$. The
per-worker bandwidth cost is asymptotically independent of $P$,
so the scalability is ultimately limited only by the latency
constant.
\end{corollary}

\begin{lemma}[ZeRO-3 communication volume]\label{lem:zero3_bw}
ZeRO-3 incurs, per training step, a communication volume of
\begin{equation}
    V_{\text{ZeRO-3}} = \frac{2 (P - 1)}{P} \cdot M s
                     + \frac{2 (P - 1)}{P} \cdot M s
                     + \frac{2 (P - 1)}{P} \cdot M s_{\text{grad}}
\end{equation}
(two all-gathers for forward/backward parameters plus one reduce-scatter
for gradients).
\end{lemma}

\begin{proof}
Every parameter is sharded across $P$ workers. Before the
forward pass, each worker must receive the other workers' shards
via an all-gather, which contributes $(P-1)/P \cdot M s$ per
worker to the communication budget. By symmetry, the analogous
all-gather issued before the backward pass contributes another
$(P-1)/P \cdot M s$. The reduce-scatter over gradients is the
dual operation and contributes
$(P-1)/P \cdot M s_{\text{grad}}$. Adding the three terms yields
the total claimed in the lemma.
\end{proof}

\begin{remark}
Comparing Lemma~\ref{lem:zero3_bw} against Lemma~\ref{lem:ddp_bw},
ZeRO-3 sends about $1.5\times$ as much data per step as DDP.
That overhead is paid off by the $8\times$ drop in memory, which
in turn permits a correspondingly larger batch size and therefore
fewer steps per epoch.
\end{remark}

\section{Exponential Decay Assumption}\label{app:proof_decay}

The assumption that attention weights decay as
$A_{ij} = A_0 e^{-\lambda |i - j|}$ is a modelling choice. It can
be justified from two directions, which we take up in turn.

\begin{proposition}[Empirical decay]\label{prop:decay}
For a transformer with softmax attention and inputs drawn from a
stationary Markov process with mixing time $\tau$, the expected
attention weight between tokens at distance $\Delta = |i - j|$
satisfies
$\mathbb{E}[P_{ij}] \le C \cdot e^{-\Delta / \tau}$ for some constant
$C > 0$.
\end{proposition}

\begin{proof}[Proof sketch]
The data-processing inequality applied to the Markov chain
implies that the mutual information $I(x_i; x_j)$ decays
exponentially in $|i - j|$ with rate $1/\tau$. Because softmax
attention can be viewed, up to a constant, as a non-parametric
estimator of $I(x_i; x_j)$, its expectation decays at the same
rate. Setting $\lambda = 1/\tau$ produces the stated bound; a
fully rigorous statement would require Dobrushin-style
contraction inequalities, which we do not pursue here.
\end{proof}

\begin{remark}
The closed-form optimal-window formula
\eqref{eq:opt_w_app} is comparatively robust to misspecification
of $\lambda$. The logarithm flattens parameter errors, so
halving or doubling $\lambda$ shifts $w^\star$ by at most
$\ln 2 \cdot 2/\lambda$---typically less than $20\%$ of the
value itself.
\end{remark}

\section{Reproducibility of the Numerical
Accuracy Bound}\label{app:reproducibility}

\begin{proposition}[FP32 error bound for SWFA]\label{prop:fp32_err}
Let $u$ denote the FP32 unit round-off ($u \approx 6 \times 10^{-8}$).
The SWFA output $\hat O$ satisfies
\begin{equation}
    \|O - \hat O\|_\infty
    \le N (3 d + 2) u \cdot
        \|Q\|_\infty \|K\|_\infty \|V\|_\infty
        + \mathcal{O}(u^2).
\end{equation}
\end{proposition}

\begin{proof}[Proof sketch]
Every entry of $S_{rc} = Q_r K_c^\top / \sqrt{d}$ is the inner
product of two $d$-vectors and therefore carries a relative
error bounded by $d u$. The online softmax introduces at most
two multiplications and one subtraction per block, adding $3u$
per block pass; summed over $T_c$ column blocks, and using
sub-additivity of rounding, the total is
$3d + 2 T_c u \le 3d + 2 N/B_c u$. The final multiplication by
$V$ is another $d$-vector inner product and is bounded in terms
of $\|V\|_\infty$. Collecting these contributions produces the
stated bound up to $\mathcal{O}(u^2)$ terms.
\end{proof}

For $N = 128$ and $d = 16$,
Proposition~\ref{prop:fp32_err} predicts a worst-case FP32 error
of roughly $4 \times 10^{-5}$, which comfortably dominates the
$\sim 8 \times 10^{-7}$ actually observed in
Table~\ref{tab:swfa_accuracy}. The gap is the usual one: the
bound assumes worst-case random inputs, which pessimistically
overestimates the error on the tame inputs used in practice.
